% Independent domain proof for the direct two-coordinate field gates.
\subsection{Direct field quotients and their domains}
\label{subsec:direct-field-quotients}

Put $k=\mathbb Q(\sqrt2)$, $d=3-\sqrt2$ and $D=d^2$.
For $a=a_0+a_1\sqrt2$ and $b=b_0+b_1\sqrt2$, multiplication
can be represented with two output coordinates and two quadratic equations:
\[
 q_0=a_0b_0+2a_1b_1,\qquad q_1=a_0b_1+a_1b_0.
\]
Similarly the inverse of $b$ is specified by
\[
 b_0q_0+2b_1q_1=1,\qquad b_0q_1+b_1q_0=0.
\]
These inverse equations themselves force $b\ne0$.  Because $k$ is
a field, they have the unique solution $q=b^{-1}$ when $b\ne0$.
If other equations already imply $b\ne0$, a quotient $q=a/b$ can
instead be specified directly by
\[
 b_0q_0+2b_1q_1=a_0,\qquad
 b_0q_1+b_1q_0=a_1.
\]
This uses only the two quotient coordinates.  A nonzero-domain argument
is essential here: without it the equations could permit $a=b=0$.

\begin{lemma}\label{lem:doubling-quotient-domain}
On $E:y^2=x^3-Dx$, the equation
\[
                           (2y)m=3x^2-D
\]
forces $y\ne0$ at every affine point, and hence defines the usual
doubling slope.  Every finite point of $2E(k)$ has nonzero abscissa
and nonzero ordinate.
\end{lemma}
\begin{proof}
If $y=0$, the curve equation gives either $x=0$ or $x^2=D$.
The numerator $3x^2-D$ is then respectively $-D$ or $2D$,
both nonzero.  Thus the slope equation cannot be satisfied when
$y=0$.  This argument uses the curve equation, and does not need
a separate assertion that $3x^2-D$ never vanishes elsewhere.

Let $Q=2R$ be finite.  The point $R$ is affine with nonzero
ordinate, and its doubling formula gives
\[
                  x(Q)=\left(\frac{x(R)^2+D}{2y(R)}\right)^2.
\]
Since $D=d^2$ and $k$ embeds in $\mathbb R$, the numerator is
nonzero.  Thus $x(Q)$ is a nonzero square in $k$.  A point with
zero ordinate has abscissa $0,d$, or $-d$.  The last two cannot
be squares in $k$: their norm to $\mathbb Q$ is $7$, whereas
the norm of a square in $k$ is a rational square.  Consequently
$y(Q)\ne0$ as well.
\end{proof}

The lemma shows that a halfpoint on the curve followed by the direct
doubling-slope equations defines a finite point of $2E(k)$ with the
required nonzero coordinates.  Both usual coordinates of the doubled
point are then given by their standard polynomial expressions in the
slope, so they lie on $E$.  Repeating this operation is valid by the
same argument.  If the usual subtraction slope is defined by its
existing abscissa-inequality guard, the resulting finite difference
of two points in $2E(k)$ also belongs to $2E(k)$ and has the same
nonvanishing properties.

Therefore the parameter $z(Q)=-x(Q)/y(Q)$ may be specified by a
direct two-coordinate quotient for every such point $Q$.  It is
itself nonzero.  For the reciprocal logarithms used in the source,
\[
              \ell_J(Q)=z(Q)G_J\bigl(D/x(Q)^2\bigr),
\]
the separately proved global nonvanishing of $G_J$ for $J\ge2$
makes $\ell_J(Q)\ne0$.  All truncations used here have
$J\in\{3,12,15\}$.  Ratios with these logarithms as denominators
may likewise use direct quotient equations.  This last step depends
on the source's stated nonvanishing lemma, not on numerical evaluation.

The abscissa inequality for the subtraction slope is retained: if two
points have the same abscissa, an unguarded quotient equation could have
both sides zero.  None of the preceding arguments removes that guard.

\subsection{Square scales without zero branches}
\label{subsec:square-scales-without-branches}

For a finite $Q\in2E(k)$ write $x(Q)=x_0+x_1\sqrt2$, with
$x_0,x_1\in\mathbb Q$.  The preceding lemma gives
$x(Q)=(u_0+u_1\sqrt2)^2$ for some nonzero element of $k$.
Consequently
\[
                    x_0=u_0^2+2u_1^2>0.
\]
The two scales
\[
                    A_0=x_0^2,\qquad A_1=(1+x_1^2)^2
\]
are therefore nonzero rational squares for every allowed point.  They
need no finite choice recording whether a coordinate is zero.

For completeness we verify the height estimate needed when using these
scales.  Let $h$ denote the logarithmic Weil height.  For rational
$v=a/b$ in lowest terms with $b>0$, the numerator and denominator
of $1+v^2=(a^2+b^2)/b^2$ are coprime.  Hence
\[
         h(1+v^2)=\log(a^2+b^2)\ge2h(v).
\]
It follows that $h(A_0)=2h(x_0)$ and $h(A_1)\ge4h(x_1)$.

Suppose that the contact estimates for a nonzero slope $\lambda$ and
both scaled slopes give, for $H,N\ge1$ and $\alpha\ge0$,
\[
 h(\lambda),\ h(A_i\lambda)\le HN^\alpha+\log3
                                      \qquad(i=0,1).
\]
Since $h(A_i)\le h(A_i\lambda)+h(\lambda)$, these imply
\[
 h(x_0)\le(H+\log3)N^\alpha,\qquad
 h(x_1)\le\tfrac12(H+\log3)N^\alpha.
\]
The standard height inequalities for sum and product give
\[
 \begin{aligned}
 h\bigl(x_0+x_1\sqrt2\bigr)
 &\le h(x_0)+h(x_1)+\tfrac32\log2\\
 &\le\tfrac32(H+\log3)N^\alpha+\tfrac32\log2
   <7HN^\alpha.
 \end{aligned}
\]
For the last inequality it suffices to use $H,N^\alpha\ge1$ and
$\log6<2$.  Thus these scales preserve the original height constant
with room to spare.  All other detector and completeness arguments use
only that each chosen scale is a fixed nonzero rational square once the
root points are fixed; that condition has been retained.  For four root
points this removes all eight zero/nonzero scale choices.
